\SourcePage{266}{269}
\section{The finite-rotation operator}

We return to spinor transformations and show how their coefficients can be expressed in terms of the angles through which the coordinate axes are rotated.

By the definition of the angular-momentum operator (here, spin), $1+i\delta\varphi\,\mathbf n\hat{\mathbf s}$ is the operator for a rotation through $\delta\varphi$ about the direction of the unit vector $\mathbf n$. For a spin-$1/2$ wave function, a spinor of rank 1, one sets $\hat{\mathbf s}=\hat{\boldsymbol\sigma}/2$. A finite rotation through $\varphi$ about the same direction is therefore represented by
\begin{equation}
 \hat U_{\mathbf n}=\exp(i\varphi\,\mathbf n\hat{\boldsymbol\sigma}/2).
 \tag{58.1}
\end{equation}
(Compare (15.13).) Like every function of the Pauli matrices (see Problem 1, §~55), this reduces to an expression linear in them:
\begin{equation}
 \hat U_{\mathbf n}=\cos(\varphi/2)+i\mathbf n\hat{\boldsymbol\sigma}\sin(\varphi/2).
 \tag{58.2}
\end{equation}
For a rotation about $z$,
\begin{equation}
 \hat U_z(\varphi)=\cos\frac\varphi2+i\hat\sigma_z\sin\frac\varphi2
 =\begin{pmatrix}e^{i\varphi/2}&0\\0&e^{-i\varphi/2}\end{pmatrix}.
 \tag{58.3}
\end{equation}
Thus the spinor components transform as
\[
 \psi^{1'}=\psi^1e^{i\varphi/2},\qquad \psi^{2'}=\psi^2e^{-i\varphi/2}.
\]
In particular, a rotation through $2\pi$ changes the sign of the spinor components; every spinor of odd rank consequently has the same property (compare the end of §~55).

Similarly, rotations through $\varphi$ about $x$ or $y$ have the matrices
\begin{equation}
 \hat U_x(\varphi)=\begin{pmatrix}\cos(\varphi/2)&i\sin(\varphi/2)\\i\sin(\varphi/2)&\cos(\varphi/2)\end{pmatrix},\qquad
 \hat U_y(\varphi)=\begin{pmatrix}\cos(\varphi/2)&\sin(\varphi/2)\\-\sin(\varphi/2)&\cos(\varphi/2)\end{pmatrix}.
 \tag{58.4}
\end{equation}

\SourcePage{267}{270}
For the special case of a rotation through $\pi$ about $y$,
\begin{equation}
 \psi^{1'}=\psi^2,\qquad \psi^{2'}=-\psi^1,
 \qquad\text{i.e.}\qquad \psi^{\lambda'}=\psi_\lambda.
 \tag{58.5}
\end{equation}

We can now write the transformation matrix for an arbitrary rotation in terms of its Euler angles.

A rotation of the axes specified by Euler angles $\alpha,\beta,\gamma$ is performed in three stages: (1) rotation through $\alpha$ ($0\leqslant\alpha\leqslant2\pi$) about $z$; (2) rotation through $\beta$ ($0\leqslant\beta\leqslant\pi$) about the new position of $y$ ($ON$ in Fig.~20, the \emph{line of nodes}); and (3) rotation through $\gamma$ ($0\leqslant\gamma\leqslant2\pi$) about the final position $z'$ of the $z$ axis\footnote{The systems $xyz$ and $x'y'z'$ are, as always, right-handed; positive angles follow the direction of a right-handed screw advancing along the positive rotation axis.

This definition of the Euler angles, customary in quantum-mechanical applications, differs from that in §~35 (Volume I) because the second rotation is about $y$, not $x$. The angles $\alpha,\beta,\gamma$ are related to the angles $\varphi,\theta,\psi$ of Volume I (not to be confused with the spherical angles $\varphi,\theta$) by $\psi=\gamma-\pi/2$.}.

\begin{figure}[ht]
\centering
\begin{tikzpicture}[scale=2.7,line cap=round,line join=round,
  ax/.style ={-{Stealth[length=2.8mm]},line width=0.75pt},
  bax/.style={-{Stealth[length=2.8mm]},line width=0.75pt,dashed},
  ang/.style={-{Stealth[length=2.1mm]},line width=0.6pt},
  ring/.style={line width=0.9pt},
  every node/.style={font=\normalsize}]
\def\al{55}
\def\be{34}
\def\ga{18}
\def\uyx{-0.36}\def\uyy{-0.66}
\def\R{1.30}
\pgfmathsetmacro\nx{sin(\al) + \uyx*cos(\al)}
\pgfmathsetmacro\ny{\uyy*cos(\al)}
\pgfmathsetmacro\mx{cos(\be)*cos(\al) - \uyx*cos(\be)*sin(\al)}
\pgfmathsetmacro\my{-\uyy*cos(\be)*sin(\al) + sin(\be)}
\pgfmathsetmacro\zpx{-sin(\be)*cos(\al) + \uyx*sin(\be)*sin(\al)}
\pgfmathsetmacro\zpy{\uyy*sin(\be)*sin(\al) + cos(\be)}
\pgfmathsetmacro\ypx{cos(\ga)*\nx + sin(\ga)*\mx}
\pgfmathsetmacro\ypy{cos(\ga)*\ny + sin(\ga)*\my}
\pgfmathsetmacro\xpx{sin(\ga)*\nx - cos(\ga)*\mx}
\pgfmathsetmacro\xpy{sin(\ga)*\ny - cos(\ga)*\my}
\draw[ring,fixedC] plot[domain=\al-84:\al+270,samples=200,smooth]
  ({\R*(-cos(\x) + \uyx*sin(\x))}, {\R*\uyy*sin(\x)});
\draw[ring,bodyC] plot[domain=240-\ga:270-\ga,samples=30,smooth]
  ({\R*(cos(\x)*\xpx + sin(\x)*\ypx)}, {\R*(cos(\x)*\xpy + sin(\x)*\ypy)});
\draw[ring,dashed,bodyC] plot[domain=270-\ga:600-\ga,samples=220,smooth]
  ({\R*(cos(\x)*\xpx + sin(\x)*\ypx)}, {\R*(cos(\x)*\xpy + sin(\x)*\ypy)});
\draw[fixedC,line width=0.75pt] (0,0) -- (1.85,0);
\draw[ax,fixedC] (0,0) -- (-1.85,0) node[left=2pt] {$X$};
\draw[ax,fixedC] (0,0) -- ({1.75*\uyx},{1.75*\uyy}) node[below=2pt] {$Y$};
\draw[ax,fixedC] (0,0) -- (0,1.90) node[above=1pt] {$Z$};
\node[right=2pt,fixedC] at (0.03,1.73) {$\boldsymbol{\delta\alpha}$};
\draw[bax,bodyC] (0,0) -- ({1.60*\xpx},{1.60*\xpy}) node[below=2pt] {$x'$};
\draw[bax,bodyC] (0,0) -- ({1.85*\ypx},{1.85*\ypy}) node[right=2pt] {$y'$};
\draw[dash dot,bodyC,line width=0.7pt] (0,0) -- (-0.6381,0.6935);
\draw[bax,bodyC] (-0.6381,0.6935) -- ({2.10*\zpx},{2.10*\zpy}) node[above left=0pt] {$z'$};
\node[bodyC] at ($({2.10*\zpx},{2.10*\zpy}) + (0.32,-0.05)$) {$\boldsymbol{\delta\gamma}$};
\draw[nodeC,-{Stealth[length=2.6mm]},line width=0.7pt]
  ({-\R*\nx},{-\R*\ny}) -- ({1.22*\R*\nx},{1.22*\R*\ny});
\node[nodeC] at ($({1.22*\R*\nx},{1.22*\R*\ny}) + (0.26,0.06)$) {$\boldsymbol{\delta\beta}$};
\node[nodeC] at ($({1.22*\R*\nx},{1.22*\R*\ny}) + (0.02,-0.17)$) {$N$};
\draw[ang] plot[domain=0:\al,samples=40,smooth]
  ({0.50*(sin(\x) + \uyx*cos(\x))}, {0.50*\uyy*cos(\x)});
\node at ($({0.50*(sin(\al/2) + \uyx*cos(\al/2))}, {0.50*\uyy*cos(\al/2)})
          + (-0.02,-0.13)$) {$\alpha$};
\draw[ang] plot[domain=0:\ga,samples=30,smooth]
  ({0.60*(cos(\x)*\nx + sin(\x)*\mx)}, {0.60*(cos(\x)*\ny + sin(\x)*\my)});
\node at ({0.76*(cos(\ga/2)*\nx + sin(\ga/2)*\mx)},
          {0.76*(cos(\ga/2)*\ny + sin(\ga/2)*\my)}) {$\gamma$};
\draw[line width=0.6pt] plot[domain=0:\be,samples=30,smooth]
  ({0.62*sin(\x)*(-cos(\al) + \uyx*sin(\al))}, {0.62*(\uyy*sin(\x)*sin(\al) + cos(\x))});
\node at ($({0.62*sin(\be/2)*(-cos(\al) + \uyx*sin(\al))},
           {0.62*(\uyy*sin(\be/2)*sin(\al) + cos(\be/2))}) + (0.06,0.11)$) {$\beta$};
\draw[bodyC,line width=0.55pt] (-0.5635,0.6126) -- (-0.6354,0.5294) -- (-0.7100,0.6103);
\node at (-0.13,-0.14) {$O$};
\end{tikzpicture}
\par\smallskip {\small Fig.~20. Definition of the Euler angles.}
\end{figure}

The angles $\alpha,\beta$ evidently coincide with the spherical angles $\varphi,\theta$ of the new $z'$ axis relative to $xyz$: $\alpha=\varphi$, $\beta=\theta$.

For this sequence of rotations, the complete transformation matrix is
\[
 \hat U(\alpha,\beta,\gamma)=\hat U_z(\gamma)\hat U_y(\beta)\hat U_z(\alpha).
\]
Direct multiplication gives
\begin{equation}
 \hat U(\alpha,\beta,\gamma)=
 \begin{pmatrix}
 \cos(\beta/2)e^{i(\alpha+\gamma)/2}&\sin(\beta/2)e^{-i(\alpha-\gamma)/2}\\
 -\sin(\beta/2)e^{i(\alpha-\gamma)/2}&\cos(\beta/2)e^{-i(\alpha+\gamma)/2}
 \end{pmatrix}.
 \tag{58.6}
\end{equation}
By definition, higher-rank spinors transform as products of rank-1 spinor components. In physical applications, however, the transformation laws of the corresponding wave functions $\psi_{jm}$ are of greater interest than those of the spinors themselves.

\SourcePage{268}{271}
Let $\psi_{jm}$, $m=j,j-1,\ldots,-j$, describe in the system $xyz$ a state with specified angular momentum $j$, and let $\psi_{jm'}$ describe the same state relative to $x'y'z'$. Thus $m$ is the value of $J_z$ and $m'$ that of $J_{z'}$. The two sets are related linearly; write
\begin{equation}
 \psi_{jm}=\sum_{m'}D^{(j)}_{m'm}\psi_{jm'}.
 \tag{58.7}
\end{equation}
The coefficients $D^{(j)}_{m'm}$ form, in the indices $m',m$, a matrix of order $2j+1$: the \emph{finite-rotation matrix} $\hat D^{(j)}$. Its elements depend on the rotation angles $\alpha,\beta,\gamma$ of $x'y'z'$ relative to $xyz$.

The finite-rotation matrix can be constructed using the spinor representation of $\psi_{jm}$.

For $j=1/2$, the two functions $\psi_{1/2,m}$ form a covariant spinor of rank 1. By (56.13), its transformation from $x'y'z'$ to $xyz$ is effected by $\hat U$ in (58.6), so $\hat D^{(1/2)}=\hat U$\footnote{The matrix indices in (58.7) are ordered precisely as required for multiplying the columns of $\hat D^{(j)}$ by the functions $\psi_{jm'}$ arranged in a row. Symbolically, (58.7) would read $\psi_{jm}=(\boldsymbol\psi_j\hat D^{(j)})_m$, consistently with (56.13).}. Write its elements as
\[
 D^{(1/2)}_{m'm}(\alpha,\beta,\gamma)=e^{im'\gamma}d^{(1/2)}_{m'm}(\beta)e^{im\alpha},
\]
where
\begin{equation}
 \begin{array}{c|cc}
 d^{(1/2)}_{m'm}(\beta)&\frac12&-\frac12\\ \hline
 \frac12&\cos(\beta/2)&\sin(\beta/2)\\
 -\frac12&-\sin(\beta/2)&\cos(\beta/2)
 \end{array}.
 \tag{58.8}
\end{equation}

For arbitrary $j$, (57.6) relates $\psi_{jm}$ to a symmetric covariant spinor of rank $2j$. Its transformation matrix is a product of $2j$ matrices $\hat D^{(1/2)}$, one acting on each spinor index. Multiplying and returning to $\psi_{jm}$ gives

\SourcePage{269}{272}
\begin{equation}
 D^{(j)}_{m'm}(\alpha,\beta,\gamma)=e^{im'\gamma}d^{(j)}_{m'm}(\beta)e^{im\alpha},
 \tag{58.9}
\end{equation}
where\footnote{The calculation is given in A.~R.~Edmonds, \emph{Angular Momentum in Quantum Mechanics}, Princeton, 1957 (see also the Russian translation of his article in \emph{Deformation of Atomic Nuclei}, Moscow, 1958). The definition of $D^{(j)}_{m'm}$ in (58.9), (58.10) differs from Edmonds's book by interchanging $\alpha$ and $\gamma$ (the more natural convention here), and from the article also by reversing the signs of all angles.}
\begin{equation}
 \begin{aligned}
 d^{(j)}_{m'm}(\beta)&=
 \left[\frac{(j+m')!(j-m')!}{(j+m)!(j-m)!}\right]^{1/2}
 \left(\cos\frac\beta2\right)^{m'+m}\\
 &\quad\times\left(\sin\frac\beta2\right)^{m'-m}
 P_{j-m'}^{(m'-m,m'+m)}(\cos\beta),
 \end{aligned}
 \tag{58.10}
\end{equation}
and
\begin{equation}
 \begin{aligned}
 P_n^{(a,b)}(\cos\beta)&=\frac{(-1)^n}{2^n n!}
 (1-\cos\beta)^{-a}(1+\cos\beta)^{-b}\\
 &\quad\times\left(\frac d{d\cos\beta}\right)^n
 [(1-\cos\beta)^{a+n}(1+\cos\beta)^{b+n}]
 \end{aligned}
 \tag{58.11}
\end{equation}
are the Jacobi polynomials\footnote{For their relation to the hypergeometric series, see §~e, formula (e.11).}. Notice that
\begin{equation}
 d^{(j)}_{m'm}(\beta)=0\quad\text{if }|m|>j\text{ or }|m'|>j.
 \tag{58.12}
\end{equation}

The functions $d^{(j)}_{m'm}$ have several symmetry properties. Although these can be read from (58.11), (58.12), it is simpler to derive them directly from their definition as rotation coefficients.

The rotation matrix $\hat D^{(j)}$ is unitary. The inverse of the rotation $(\alpha,\beta,\gamma)$ is $(-\gamma,-\beta,-\alpha)$; since $\hat d$ is real, this gives
\begin{equation}
 d^{(j)}_{m'm}(-\beta)=d^{(j)}_{mm'}(\beta).
 \tag{58.13}
\end{equation}

\SourcePage{270}{273}
We also have
\begin{equation}
 d^{(j)}_{m'm}(\beta)=d^{(j)}_{-m,-m'}(\beta),
 \tag{58.14}
\end{equation}
and
\begin{equation}
 d^{(j)}_{m'm}(\pi)=(-1)^{j+m}\delta_{m',-m},\qquad
 d^{(j)}_{m'm}(-\pi)=(-1)^{j-m}\delta_{m',-m},\qquad
 d^{(j)}_{m'm}(0)=\delta_{m'm}.
 \tag{58.15}
\end{equation}
For $j=1/2$ these follow directly from (58.8), and for arbitrary $j$ from the construction of the transformation matrix given above.

Perform a rotation through $\pi-\beta$ as successive rotations through $\pi$ and $-\beta$:
\[
 d^{(j)}_{m'm}(\pi-\beta)=\sum_{m''}d^{(j)}_{m'm''}(\pi)d^{(j)}_{m''m}(-\beta)
 =(-1)^{j-m'}d^{(j)}_{-m',m}(-\beta),
\]
or, using (58.13),
\begin{equation}
 d^{(j)}_{m'm}(\pi-\beta)=(-1)^{j-m'}d^{(j)}_{m,-m'}(\beta).
 \tag{58.16}
\end{equation}
Two rotations about the same axis give the same result in either order. Reversing the order of the rotations $-\beta$ and $\pi$ and comparing with (58.16), we obtain
\begin{equation}
 d^{(j)}_{m'm}(\beta)=(-1)^{m'-m}d^{(j)}_{-m',-m}(\beta).
 \tag{58.17}
\end{equation}
Equations (58.17), (58.14), and (58.13) imply
\begin{equation}
 d^{(j)}_{m'm}(\beta)=(-1)^{m'-m}d^{(j)}_{mm'}(\beta)
 =(-1)^{m'-m}d^{(j)}_{m'm}(-\beta).
 \tag{58.18}
\end{equation}
These relations give corresponding symmetries of the complete functions $D^{(j)}_{m'm}$. In particular,
\begin{equation}
 D^{(j)*}_{m'm}(\alpha,\beta,\gamma)=D^{(j)}_{m'm}(-\alpha,\beta,-\gamma)
 =(-1)^{m'-m}D^{(j)}_{-m',-m}(\alpha,\beta,\gamma).
 \tag{58.19}
\end{equation}
Mathematically, the matrices $\hat D^{(j)}$ furnish unitary irreducible representations of the rotation group of dimension $2j+1$ (see §~98). Hence the orthogonality and normalization relation is
\begin{equation}
 \int D^{(j_1)*}_{m'_1m_1}(\alpha,\beta,\gamma)
 D^{(j_2)}_{m'_2m_2}(\alpha,\beta,\gamma)\frac{d\omega}{8\pi^2}
 =\frac1{2j_1+1}\delta_{j_1j_2}\delta_{m_1m_2}\delta_{m'_1m'_2},
 \tag{58.20}
\end{equation}
where $d\omega=\sin\beta\,d\alpha\,d\beta\,d\gamma$.

\SourcePage{271}{274}
Orthogonality in $m$ and $m'$ is supplied by the factor $\exp\{i(m\alpha+m'\gamma)\}$. Orthogonality in $j$ belongs to the functions $d^{(j)}_{m'm}$:
\begin{equation}
 \int_0^\pi d^{(j_1)}_{m'm}(\beta)d^{(j_2)}_{m'm}(\beta)
 \sin\beta\,d\beta=\frac2{2j_1+1}\delta_{j_1j_2}.
 \tag{58.21}
\end{equation}

For reference, we finally list the functions $d^{(j)}_{m'm}$ for several special values of their parameters. For $j=1$,
\begin{equation}
\begin{array}{c|ccc}
 d^{(1)}_{m'm}(\beta)&1&0&-1\\ \hline
 1&\frac12(1+\cos\beta)&\frac1{\sqrt2}\sin\beta&\frac12(1-\cos\beta)\\
 0&-\frac1{\sqrt2}\sin\beta&\cos\beta&\frac1{\sqrt2}\sin\beta\\
 -1&\frac12(1-\cos\beta)&-\frac1{\sqrt2}\sin\beta&\frac12(1+\cos\beta)
\end{array}
\tag{58.22}
\end{equation}
For integral $j=l$ and $m'=0$, (58.10), (58.11) give
\begin{equation}
 d^{(l)}_{0m}(\beta)=\left[\frac{(l-m)!}{(l+m)!}\right]^{1/2}P_l^m(\cos\beta).
 \tag{58.23}
\end{equation}
The origin of this formula is immediate from (58.7). Refer the functions $\psi_{jm'}$ on its right-hand side to the $z'$ axis. For $j=l$,
\begin{equation}
 \psi_{lm'}(\mathbf n_{z'})=i^l\sqrt{\frac{2l+1}{4\pi}}\delta_{m'0}.
 \tag{58.24}
\end{equation}
The function $\psi_{jm}$ on the left is then the spherical harmonic $Y_{lm}(\beta,\alpha)$ of the spherical angles $\varphi=\alpha$, $\theta=\beta$ of $z'$. Substitution of (58.24) in (58.7) gives
\begin{equation}
 Y_{lm}(\beta,\alpha)=i^l\sqrt{\frac{2l+1}{4\pi}}D^{(l)}_{0m}(\alpha,\beta,\gamma),
 \tag{58.25}
\end{equation}
which is equivalent to (58.23).

\SourcePage{272}{275}
Finally, when one index $m$ or $m'$ has its largest possible value,
\begin{equation}
 d^{(j)}_{jm}(\beta)=(-1)^{j-m}d^{(j)}_{-j,m}(\beta)
 =\left[\frac{(2j)!}{(j+m)!(j-m)!}\right]^{1/2}
 \left(\cos\frac\beta2\right)^{j+m}
 \left(\sin\frac\beta2\right)^{j-m}.
 \tag{58.26}
\end{equation}
